\documentclass[10pt,a4paper]{amsart}
\usepackage[margin=19mm]{geometry}
\usepackage{amsmath,amssymb,mathtools}
\usepackage[scr=rsfs,cal=euler]{mathalfa}
\usepackage{xfrac}
\usepackage[unicode,hidelinks,bookmarks=false]{hyperref}
\newcommand{\ie}{i.e.\ }
\newcommand{\cf}{cf.\ }
\newcommand{\resp}{respectively\ }
\newcommand{\Subsection}{\subsection}
\providecommand{\nobreakdash}{\nobreak\hbox{-}}
\setlength{\emergencystretch}{2em}
\multlinegap0pt
\usepackage{bm}
\usepackage{bbm}
\usepackage{dsfont}
\usepackage{xspace}
\usepackage{cancel}
\usepackage{mathtools}
\usepackage[shortlabels]{enumitem}
\usepackage{amsthm}
\makeatletter
\@ifundefined{theorem}{\newtheorem{theorem}{Theorem}}{}
\@ifundefined{lemma}{\newtheorem{lemma}[theorem]{Lemma}}{}
\makeatother
\title[Rotationally symmetric translating solitons of fully nonline]{Rotationally symmetric translating solitons of fully nonlinear extrinsic geometric flows: Classification and Applications}
\author{Jos\'e Torres Santaella}
\date{Source: arXiv:2512.23623v3 (CC BY 4.0)}
\begin{document}
\maketitle

\noindent\emph{Source and licence.} Rotationally symmetric translating solitons of fully nonlinear extrinsic geometric flows: Classification and Applications. Version arXiv:2512.23623v3 (CC BY 4.0), distributed under the Creative Commons Attribution 4.0 International licence, as recorded in the arXiv licence metadata for this exact version and shown on the abstract page https://arxiv.org/abs/2512.23623v3. Full rights notice: licenses/arxiv-2512.23623-CC-BY-4.0.txt.\par\smallskip

\noindent\textit{Editorial scope.} Counted unit: the Lemma whose source label is \texttt{lem:start-moving-plane-at-infinity} in the article (source0.tex lines 4635--4666, source label \texttt{lem:start-moving-plane-at-infinity}), delivered together with the article's own complete original proof of it (source0.tex lines 4668--4804, one \texttt{proof} environment of 137 lines). No mathematics is written, repaired or re-derived: statement and proof are the article's own text line for line. Required context retained uncounted: eq:C2-asymptotic-to-bowl (lines 4603--4612). Every definition, display and label of those lines is retained unchanged. Omitted from this package and not redistributed: 1--4602, 4613--4634, 4667--4667, 4805--5845. Nothing inside a retained excerpt is omitted or altered: every retained body is delivered exactly as the article's lines are written, so no omission receipt of any kind accompanies this package. Citations: the retained excerpts carry no citation command at all, so no bibliography entry of the article is transcribed and no reference is redistributed. Reference adaptation: none was needed and none was made; every reference inside the delivered unit resolves inside it, so no unresolved reference or citation is produced. Printed slips kept literal: no printed word, symbol, hypothesis or number of the counted statement or of its proof was altered. Layout and class: the article's own class is \texttt{amsart} with packages including fontenc, inputenc, babel, amsthm, amssymb, verbatim, mathtools, mathrsfs, enumitem, graphicx, subcaption, xcolor, tikz, pgfplots; the wrapper is the lane's pinned \texttt{amsart} editorial wrapper at 10pt on A4, which restates the theorem-like environments the retained text needs and adds the article's own macro definitions that the retained text needs (none), with extra wrapper packages (bm, bbm, dsfont, xspace, cancel, mathtools, enumitem). The delivered numbering is the unit's own. Assets: no retained span contains a figure environment, a graphics-inclusion command, a PDF-page inclusion, a table or a TikZ picture; no image file is copied into this package, the package declares no asset dependency and no image file is redistributed with it. Licence: version arXiv:2512.23623v3 of this article is distributed under the Creative Commons Attribution 4.0 International licence, as recorded in the arXiv licence metadata for this exact version and shown on the abstract page https://arxiv.org/abs/2512.23623v3. A full rights notice is delivered at licenses/arxiv-2512.23623-CC-BY-4.0.txt. Characters: every retained excerpt is pure ASCII, so the delivered engine, XeLaTeX with its default Latin Modern text font, reports no missing character.
	\begin{align}
		\label{eq:C2-asymptotic-to-bowl}
		u(x)=u_B(|x|)+\varphi(x),
		\qquad
		\varphi(x)\to0,
		\qquad
		D\varphi(x)\to0,
		\qquad
		D^2\varphi(x)\to0
	\end{align}

	\begin{lemma}
		\label{lem:start-moving-plane-at-infinity}
		Assume that the positive end of $B$ is in one of the following regimes:
		\begin{enumerate}[label=\textup{(\roman*)}]
			\item the normalized $1$-nondegenerate case with $\alpha\ge1$ and
			\begin{align*}
				v_B(r)=r^\alpha-\frac{a}{r^\alpha}+\frac{b}{r^{3\alpha}}+o(r^{-4\alpha});
			\end{align*}
			
			\item the $1$-degenerate power-growth case with
			\begin{align*}
				v_B(r)
				=
				A_\gamma
				r^{
					\frac{\alpha(k_\gamma+1)}
					{k_\gamma+1-2\alpha}
				}
				(1+o(1)),
				\mbox{ with }
				\frac{\alpha(k_\gamma+1)}
				{k_\gamma+1-2\alpha}\ge1;
			\end{align*}
			
			\item the critical $1$-degenerate case with
			$\log v_B(r)=\frac{c_\gamma}{2\alpha}r^{2\alpha}+o(r^{2\alpha})$.
		\end{enumerate}
		Then there exists $t_0>0$ such that $u_t^*(x)>u(x)$ for every
		$t\ge t_0$ and every $x_1<t$.  Moreover, for every fixed $t>0$ and every
		$\sigma\in(0,t)$, there exist $R_{\sigma,t}>0$ and $\mu_{\sigma,t}>0$ such
		that $u_t^*(x)-u(x)\ge\mu_{\sigma,t}$ whenever $|x|\ge R_{\sigma,t}$ and $x_1\le t-\sigma$.
	\end{lemma}

	\begin{proof}
		For a fixed $t>0$, set $x^t=(2t-x_1,x_2,\ldots,x_n)$,
		$\rho=|x|$, and $\rho_t=|x^t|$.  Then
		\begin{align}
			\label{eq:reflected-radius-difference}
			\rho_t^2-\rho^2=4t(t-x_1).
		\end{align}
		In particular, $\rho_t>\rho$ whenever $x_1<t$, and
		$u_B(\rho_t)-u_B(\rho)=\int_\rho^{\rho_t}v_B(s)\,ds$.
		
		\medskip
		
	Firstly, we prove the separation at infinity.  Let $\sigma\in(0,t)$ and assume $x_1\le t-\sigma$.  Then $\rho_t^2-\rho^2\ge4t\sigma$. In the normalized $1$-nondegenerate case with $\alpha>1$, the expansion gives $v_B(s)=s^\alpha(1+o(1))$.  Hence, for every fixed $\varepsilon\in(0,1)$ and all large $s$, $v_B(s)\ge(1-\varepsilon)s^\alpha$. For $s\in[\rho,\rho_t]$, the inequality $s^\alpha\ge\rho^{\alpha-1}s$ gives
		\begin{align*}
			u_B(\rho_t)-u_B(\rho)
			&\ge
			(1-\varepsilon)\rho^{\alpha-1}\int_\rho^{\rho_t}s\,ds  \\
			&=
			\frac{(1-\varepsilon)}{2}\rho^{\alpha-1}(\rho_t^2-\rho^2)
			\ge
			2(1-\varepsilon)t\sigma\,\rho^{\alpha-1}.
		\end{align*}
		Thus $u_B(\rho_t)-u_B(\rho)\to+\infty$ uniformly along sequences with
		$|x|\to\infty$ and $x_1\le t-\sigma$. On the other hand, if  $\alpha=1$,
		$v_B(s)=s-as^{-1}+bs^{-3}+o(s^{-4})$.  Integrating from $\rho$ to
		$\rho_t$, we get
		\begin{align*}
			u_B(\rho_t)-u_B(\rho)
			=
			\frac{\rho_t^2-\rho^2}{2}
			-
			a\log\frac{\rho_t}{\rho}
			-
			\frac{b}{2}
			\left(\frac{1}{\rho_t^2}-\frac{1}{\rho^2}\right)
			+
			o(1).
		\end{align*}
		Here $\rho_t\rho^{-1}\to1$ uniformly as $|x|\to\infty$ in
		$\{x_1\le t-\sigma\}$.  Indeed,
		$\rho_t^2-\rho^2=4t(t-x_1)\le4t(t+\rho)$, and hence
		$\rho_t^2\rho^{-2}=1+O(\rho^{-1})$.  Therefore the logarithmic and
		inverse-square terms vanish at infinity, and
		\begin{align*}
			\liminf\limits_{\substack{|x|\to\infty\\ x_1\le t-\sigma}}
			\bigl(u_B(\rho_t)-u_B(\rho)\bigr)
			\ge
			2t\sigma .
		\end{align*}
		Next, in the $1$-degenerate power-growth case, if the exponent
		$\frac{\alpha(k_\gamma+1)}{k_\gamma+1-2\alpha}$ is strictly larger than $1$,
		the preceding superlinear estimate applies with that exponent and gives
		$u_B(\rho_t)-u_B(\rho)\to+\infty$ uniformly on escaping sequences with
		$x_1\le t-\sigma$.  When the exponent is equal to $1$,
		$v_B(s)=A_\gamma s(1+o(1))$, and the same computation as in the linear case
		gives
		\begin{align*}
			\liminf\limits_{\substack{|x|\to\infty\\ x_1\le t-\sigma}}
			\bigl(u_B(\rho_t)-u_B(\rho)\bigr)
			\ge
			2A_\gamma t\sigma .
		\end{align*}
		Finally, in the critical $1$-degenerate case, choose
		$\kappa\in\left(0,\frac{c_\gamma}{2\alpha}\right)$.  The expansion of
		$\log v_B$ gives $v_B(s)\ge\exp(\kappa s^{2\alpha})$ for all large $s$.
		For all large $\rho$ with $x_1\le t-\sigma$, the estimate
		$\rho_t+\rho\le C_t\rho$ gives
		\begin{align*}
			\rho_t-\rho
			=
			\frac{\rho_t^2-\rho^2}{\rho_t+\rho}
			\ge
			\frac{4t\sigma}{C_t\rho}.
		\end{align*}
		Consequently,
		\begin{align*}
			u_B(\rho_t)-u_B(\rho)
			\ge
			(\rho_t-\rho)\exp(\kappa\rho^{2\alpha})
			\ge
			\frac{4t\sigma}{C_t\rho}\exp(\kappa\rho^{2\alpha})
			\to+\infty .
		\end{align*}
		
		In all cases, there exist $R_{\sigma,t}>0$ and $\mu_{\sigma,t}>0$ such that
		$u_B(\rho_t)-u_B(\rho)\ge2\mu_{\sigma,t}$ whenever
		$|x|\ge R_{\sigma,t}$ and $x_1\le t-\sigma$.  From
		\eqref{eq:C2-asymptotic-to-bowl}, after increasing $R_{\sigma,t}$ if
		necessary, we obtain
		\begin{align*}
			u_t^*(x)-u(x)
			=
			u_B(\rho_t)-u_B(\rho)+\varphi(x^t)-\varphi(x)
			\ge
			\mu_{\sigma,t}
		\end{align*}
		on the same region.
		
		\medskip
		
		Next, we obtain the ordering from a far plane.  The expansions above imply that
		there are constants $c>0$ and $r_0>0$ such that $v_B(s)\ge cs$ for
		$s\ge r_0$.  Hence there is $C>0$ such that, for all $0\le\rho<\rho_t$,
		\begin{align*}
			u_B(\rho_t)-u_B(\rho)
			\ge
			\frac{c}{2}(\rho_t^2-\rho^2)-C .
		\end{align*}
		For $x_1\le\frac{t}{2}$, one has $\rho_t^2-\rho^2\ge2t^2$.  The boundedness
		of $\varphi$ gives $u_t^*>u$ on $\{x_1\le\frac{t}{2}\}$ for every
		sufficiently large $t$.
		
		\medskip
		
		It remains for us to consider the strip $\frac{t}{2}<x_1<t$.  For
		$s\in[x_1,2t-x_1]$ and $x'=(x_2,\ldots,x_n)$, we have
		$\sqrt{s^2+|x'|^2}\ge s>\frac{t}{2}$.  Taking $t$ large enough so that this
		radius is at least $r_0$, the lower bound $v_B(r)\ge cr$ gives
		\begin{align*}
			v_B\left(\sqrt{s^2+|x'|^2}\right)
			\frac{s}{\sqrt{s^2+|x'|^2}}
			\ge
			cs
			\ge
			\frac{c}{2}t .
		\end{align*}
		Using $D\varphi\to0$ and increasing $t$ once more, we get
		$\partial_1u(s,x')>0$ on the strip.  Therefore
		\begin{align*}
			u_t^*(x)-u(x)
			=
			\int_{x_1}^{2t-x_1}\partial_1u(s,x')\,ds
			>
			0
		\end{align*}
		whenever $\frac{t}{2}<x_1<t$.  Combining both regions proves the lemma.
	\end{proof}

\end{document}
